\documentclass[10pt,letterpaper,oneside]{report}
\usepackage[T1]{fontenc}
\usepackage[utf8]{inputenc}
\usepackage{lmodern}
\usepackage{microtype}
\usepackage[margin=0.85in,headheight=15pt]{geometry}
\usepackage{enumitem,booktabs,longtable,tabularx,array,makecell}
\usepackage{fancyhdr,lastpage,xcolor,hyperref,titlesec}
\usepackage{listings}
\usepackage[most]{tcolorbox}
\definecolor{ink}{HTML}{182D40}
\definecolor{accent}{HTML}{126B77}
\definecolor{pale}{HTML}{EEF5F5}
\definecolor{grayblue}{HTML}{677887}
\definecolor{codebg}{HTML}{F5F7F8}
\hypersetup{colorlinks=true,linkcolor=accent,urlcolor=accent,pdfauthor={C23 Engineering Companion},pdfsubject={ISO/IEC 9899:2024 educational engineering companion}}
\titleformat{\chapter}[display]{\normalfont\huge\bfseries\color{ink}}{\filleft\large\color{accent}CHAPTER \thechapter}{1ex}{\titlerule[1pt]\vspace{1ex}}
\titleformat{\section}{\Large\bfseries\color{ink}}{\thesection}{0.7em}{}
\titleformat{\subsection}{\large\bfseries\color{accent}}{\thesubsection}{0.65em}{}
\setlist[itemize]{leftmargin=1.3em,itemsep=2pt,topsep=3pt}
\setlist[enumerate]{leftmargin=1.6em,itemsep=3pt,topsep=3pt}
\setlength{\parskip}{5pt}
\setlength{\parindent}{0pt}
\renewcommand{\arraystretch}{1.18}
\pagestyle{fancy}\fancyhf{}\fancyhead[L]{\small\sffamily C23 ENGINEERING SERIES}\fancyhead[R]{\small\sffamily\color{accent}\DocShort}\fancyfoot[L]{\small\color{grayblue}ISO/IEC 9899:2024 companion}\fancyfoot[R]{\small\thepage\ / \pageref{LastPage}}
\newtcolorbox{keybox}{breakable,colback=pale,colframe=accent,boxrule=.6pt,arc=1mm,left=8pt,right=8pt,top=6pt,bottom=6pt}
\newtcolorbox{warnbox}{breakable,colback=codebg,colframe=grayblue,boxrule=.6pt,arc=1mm,left=8pt,right=8pt,top=6pt,bottom=6pt}
\lstdefinestyle{cstyle}{language=C,basicstyle=\ttfamily\footnotesize,keywordstyle=\color{accent}\bfseries,commentstyle=\color{grayblue}\itshape,stringstyle=\color{ink},backgroundcolor=\color{codebg},frame=single,rulecolor=\color{grayblue!40},breaklines=true,showstringspaces=false,tabsize=4,columns=fullflexible,keepspaces=true,upquote=true}
\lstset{style=cstyle}
\newcommand{\refcl}[1]{\textsf{\footnotesize[ISO/IEC 9899:2024, #1]}}
\newcommand{\standard}{\textbf{[STANDARD]}}
\newcommand{\optional}{\textbf{[OPTIONAL]}}
\newcommand{\guidance}{\textbf{[GUIDANCE]}}
\newcommand{\testtag}{\textbf{[TEST]}}
\newcommand{\term}[1]{\textbf{#1}}
\newcommand{\DocShort}{CLAUSE STUDY GUIDE}
\begin{document}
\begin{titlepage}
\color{ink}\vspace*{1.2in}
{\sffamily\small\color{accent}\bfseries TECHNICAL DOCUMENTATION SUITE \hfill EDITION 2026-10}\par
\vspace{.3in}\rule{\linewidth}{1.5pt}\par\vspace{.22in}
{\Huge\bfseries C23 Clause-by-Clause Study Guide}\par\vspace{.22in}
{\Large\color{grayblue}Source-navigation cards, interpretive questions and practical assessments}\par\vspace{.4in}
\begin{keybox}\textbf{Standards basis:} ISO/IEC 9899:2024 (C23). Independent educational and engineering companion. Not an official ISO interpretation, full specification, certification or conformance assessment.\end{keybox}
\vfill{\large\bfseries C23 Software Engineering and Language Standards}\par\vspace{.15in}{\color{grayblue}Prepared October 2026 \quad | \quad Source edition: October 2024}\par
\end{titlepage}
\pagenumbering{roman}\tableofcontents\clearpage\pagenumbering{arabic}

\chapter*{Sources and interpretation key}
\addcontentsline{toc}{chapter}{Sources and interpretation key}
\textbf{Primary source:} ISO/IEC 9899:2024(en), \emph{Information technology --- Programming languages --- C}, fifth edition (2024). Clause and annex identifiers refer to the supplied licensed standard. This handbook contains original explanations, designs and examples; it does not reproduce the standard. The licensed standard remains the authoritative source for exact wording and conformance decisions.\par
\textbf{Secondary engineering sources:} \url{https://gcc.gnu.org/onlinedocs/gcc/Standards.html}; \url{https://clang.llvm.org/c_status.html}; \url{https://cmake.org/cmake/help/latest/prop_tgt/C_STANDARD.html}. Compiler support and diagnostics are release-specific: verify the installed compiler, target and library before relying on a feature.\par
\begin{keybox}\textbf{Interpretation categories.} \textbf{[STANDARD]} summarizes a requirement, constraint, semantic rule or defined facility and cites the applicable clause. \textbf{[OPTIONAL]} describes conditional or implementation-dependent support. \textbf{[GUIDANCE]} records an engineering recommendation, not an ISO mandate. \textbf{[TEST]} is an original check or exercise; a passing test is evidence but is not proof of full conformity.\end{keybox}
\textbf{Copyright boundary:} The supplied PDF is individually licensed. Do not redistribute it with this suite; distribute only these independently written companion materials in accordance with applicable rights.\par


\chapter{How to study the source}
\section{Read clauses in a consistent sequence}
For each subclause, first identify \emph{syntax}, then \emph{constraints}, then \emph{semantics}, then any references, examples, footnotes and implementation choices. For library subclauses, include \emph{description}, \emph{returns} and \emph{runtime-constraints}. Treat the exact text of the source document as decisive when the overview below is shorter.
\begin{keybox}\textbf{Study record for each clause:} identifier; concept; entities; preconditions; translation diagnostics; runtime failure modes; portability choices; one positive test; one negative test; engineering policy. Label each conclusion \textbf{STANDARD}, \textbf{OPTIONAL} or \textbf{GUIDANCE}.\end{keybox}
\chapter{Clause navigation and knowledge checks}
\small
\begin{longtable}{p{.11\linewidth}p{.19\linewidth}p{.34\linewidth}p{.24\linewidth}}\toprule Clause & Focus & Demonstrate mastery & Socratic check\\\midrule\endhead
\textbf{1--4} & Scope, definitions, conformance & Distinguish implementation conformance from program conformance; classify five kinds of behavior. & Can the compiler accept an invalid program after diagnosing it?\\\hline
\textbf{5.1--5.2} & Translation, startup, execution & Trace conceptual translation phases and explain hosted versus freestanding constraints. & Must a compiler implement physical passes for each phase?\\\hline
\textbf{5.3} & Environment limits and characters & Inventory target character sets, limits and signal assumptions. & Which limits must be checked per implementation?\\\hline
\textbf{6.1--6.2} & Notation, scope, types, representations & Diagram lifetime, linkage, alignment, qualifiers and compatibility. & What changes when an object ends lifetime but a pointer remains?\\\hline
\textbf{6.3} & Conversions & Compute promoted types in mixed signed/unsigned expressions. & Is casting the result sufficient to prevent overflow?\\\hline
\textbf{6.4} & Lexical elements & Classify tokens and encoding-dependent spelling. & When does a number become a preprocessing number?\\\hline
\textbf{6.5} & Expressions & Mark side effects, sequencing and valid pointer ranges. & What is the difference between one-past and dereferenceable?\\\hline
\textbf{6.6} & Constant expressions & Separate syntactic constants, integer constant expressions and C23 additions. & Where is compile-time evaluation required?\\\hline
\textbf{6.7.1--6.7.9} & Declarations & Construct declarators, resolve types and linkage. & What does const qualify in a nested pointer declaration?\\\hline
\textbf{6.7.10} & Type inference & Evaluate scope for auto, typeof and typeof\_unqual. & When should public interfaces avoid inference?\\\hline
\textbf{6.7.11--6.7.13} & Initialization, assertions, attributes & Contrast required semantics with optional attribute support. & What is guaranteed if an attribute is unsupported?\\\hline
\textbf{6.8--6.9} & Statements and external definitions & Trace labels, blocks, functions and translation-unit boundaries. & How does an empty parameter list differ from old C?\\\hline
\textbf{6.10} & Preprocessing & Trace phase-4 macro expansion and C23 new directives. & Does a successful \#include probe prove library support?\\\hline
\textbf{6.11} & Future language directions & List usages marked for caution rather than already banned. & Is an obsolescent feature automatically prohibited?\\\hline
\textbf{7.1--7.11} & Library fundamentals & Study reserved names, errno, locale and environment. & May a library function be called outside its parameter domain?\\\hline
\textbf{7.12--7.17} & Math, nonlocal jumps, atomics & Identify math-error behavior, atomic order and jump validity. & Can volatile replace atomic synchronization?\\\hline
\textbf{7.18--7.22} & Bits, checked arithmetic, integer widths & Prototype overflow-safe arithmetic and bit operations. & How do you guard unavailable C23 headers?\\\hline
\textbf{7.23--7.27} & I/O, allocation, strings & Trace file errors, ownership, allocation and null termination. & Does a successful realloc always preserve pointer identity?\\\hline
\textbf{7.28--7.33} & Threads, time, Unicode, library evolution & Explain thread safety, conversion state and future directions. & Are all hosted libraries equally available on freestanding targets?\\\hline
\textbf{Annex A--E} & Informative summaries & Use grammar, library tables and implementation limits as navigation aids. & Can an informative annex replace normative wording?\\\hline
\textbf{Annex F--H} & Floating-point annexes (normative) & Read applicability and floating environment requirements. & Is every floating-point implementation an Annex F implementation?\\\hline
\textbf{Annex I--J} & Warnings and portability (informative) & Build a portability-risk register. & Do Annex J examples enumerate every possible UB case?\\\hline
\textbf{Annex K--L} & Optional conditional annex facilities (normative) & Differentiate conditional extensions from core obligations. & Does Annex K being normative mean all implementations provide it?\\\hline
\textbf{Annex M} & Change history (informative) & Trace each change to its operative clause before claiming conformance. & Is a change-history bullet itself sufficient for conformance?\\\hline
\bottomrule\end{longtable}\normalsize
\chapter{Detailed study cards}
\section{Concept card: diagnostic versus undefined behavior}
\standard\ Constraints and syntax requirements can demand diagnostics, even when the construct also produces undefined behavior. By contrast, many runtime errors are not diagnosed. \refcl{4, 5.2.1.3}.\par
\testtag\ Translate one file containing a constraint violation, record every diagnostic and whether the compiler still emits code; then run a separate in-bounds/out-of-bounds sanitizer demonstration. Explain why neither outcome settles the other question.
\section{Concept card: object lifetime and pointer validity}
\standard\ Differentiate an allocated region, a declared object, and the lifetime of the designated object. \refcl{6.2.4, 6.5.7, 7.24.4}.\par
\testtag\ Draw object lifetimes for local arrays, compound literals, static objects and allocated objects. Mark which pointer evaluations remain valid after a scope or allocation ends.
\section{Concept card: integer width versus integer rank}
\standard\ Widths, ranks, promotions and value ranges are related but not interchangeable. \refcl{6.2.5, 6.3.1}.\par
\testtag\ Compare \texttt{size\_t}, \texttt{ptrdiff\_t}, \texttt{uint32\_t}, \texttt{unsigned int} and \texttt{\_BitInt(17)} on two selected targets. Provide evidence of available types and conversions.
\section{Concept card: preprocessing does not mean parsing}
\standard\ Tokenization and macro expansion occur before later semantic analysis in the conceptual phases. \refcl{5.2.1.1, 6.10}.\par
\testtag\ Use the preprocessor output to explain why a side-effecting argument may be duplicated by a function-like macro. Replace it with an inline function and compare behavior.
\section{Concept card: optional attributes}
\optional\ Standard attributes do not have universal support obligations simply because the syntax is standardized. \refcl{6.7.13}.\par
\testtag\ Gate a \texttt{nodiscard} declaration behind \texttt{\_\_has\_c\_attribute}, then compare diagnostics on two compilers.
\section{Concept card: conditional normative annexes}
\optional\ Annexes F--H, K and L are identified as normative, but application and conditional support must be determined from the annex and implementation. \refcl{Annex F--H, K--L}.\par
\testtag\ Write a one-page declaration of applicability for a hosted platform and an embedded freestanding platform; include evidence and unresolved questions.
\chapter{Worked knowledge-check solutions}
\section{Why does a warning not automatically mean nonconformance?}
A warning may be a quality or extension diagnostic beyond the standard's minimum. The important questions are whether a mandated diagnostic was produced for a violation, whether valid input programs are translated consistently with the implementation's claimed conformance, and whether behavior follows the documented contracts. \refcl{4, 5.2.1.3}.
\section{What is wrong with \texttt{sizeof(ptr)} as an array length?}
The expression measures the pointer object, not the number of elements in an array. If a function parameter accepts a pointer, the element count must be passed, derivable from another contract or otherwise proven. \refcl{6.5.4, 6.7.7}.
\section{Why is a successful sanitizer run not proof of correctness?}
Dynamic tests exercise specific executions and instrumentation. Untested inputs, optimizer paths, implementation-defined behavior and unspecified choices may remain. \guidance\ Combine static analysis, compile-time diagnostics, review, differential builds and property-based tests.
\chapter{Assessments and evidence artifacts}
\section{Written examination blueprint}
\begin{itemize}\item 20 percent: terms, conformance and diagnostics.
\item 30 percent: expressions, types, lifetimes and undefined behavior.
\item 20 percent: declarations, preprocessing and C23 change history.
\item 20 percent: libraries, threading, resources, portability.
\item 10 percent: annex applicability and evidence traceability.\end{itemize}
\section{Practical examination}
Supply a short project containing type-conversion hazards, a lifetime error, a macro side effect, an unchecked allocation, one C23 facility and one portability assumption. The learner must produce: corrected patch; clause references; before/after diagnostics; functional tests; an implementation-defined behavior note; and a design rationale.
\section{Completion standard}
\guidance\ Require 85\%+ written accuracy, all critical safety defects resolved, all cited clauses independently checked, and an evidence package that can be reproduced on a clean toolchain. This threshold is a training policy, not an ISO requirement.

\end{document}
